%======================================================================
% MONITOREO
%======================================================================

\chapter{Monitoreo}

El módulo de monitoreo proporciona visibilidad profunda del estado del sistema, permitiendo a los administradores mantener el control sobre la infraestructura y responder rápidamente a cualquier anomalía.

\section{Dashboard}

El dashboard de monitoreo presenta una vista consolidada de todas las métricas relevantes del sistema:

\begin{itemize}
    \item Uso de CPU y memoria
    \item Conexiones activas
    \item Tráfico de red
    \item Estado de servicios
    \item Alertas pendientes
\end{itemize}

La monitorización proactiva permite detectar anomalías antes de que impacten a los usuarios, reduciendo el tiempo de inactividad no planificado.

\section{Métricas en Tiempo Real}

El sistema recopila y displaying métricas en tiempo real:

\begin{itemize}
    \item \textbf{Latencia}: Tiempo de respuesta de cada componente
    \item \textbf{Throughput}: Volumen de transacciones por segundo
    \item \textbf{Errores}: Tasa de errores y excepciones
    \item \textbf{Disponibilidad}: Porcentaje de tiempo operativo
\end{itemize}

\section{Heartbeat y Latencia}

\begin{itemize}
\item \textbf{Detección en Tiempo Real}: Monitoreo continuo de la latencia entre nodos del clúster y clientes.
\item \textbf{Alertas Automáticas}: Notificaciones ante degradaciones de servicio que excedan umbrales configurados.
\item \textbf{Histórico de Métricas}: Almacenamiento de datos históricos para análisis de tendencias.
\end{itemize}

La monitorización proactiva permite detectar anomalías antes de que impacten a los usuarios, reduciendo el tiempo de inactividad no planificado.

\section{Análisis de Tráfico}

El módulo de análisis de tráfico proporciona:

\begin{itemize}
\item \textbf{Visualización de Métricas}: Gráficos en tiempo real del tráfico de red.
\item \textbf{Identificación de Cuellos de Botella}: Detección de puntos de congestión en la transferencia de datos.
\item \textbf{Integración gRPC}: Métricas específicas del protocolo de alto rendimiento.
\end{itemize}

El análisis de tráfico permite optimizar el rendimiento de la red identificando patrones de uso y potenciales cuellos de botella.

\section{Salud de Nodos}

El sistema supervisa constantemente el estado de los servicios backend:

\begin{itemize}
\item \textbf{Estado de Sandra Server}: Disponibilidad y rendimiento del middleware.
\item \textbf{Balanceo de Carga}: Verificación de que las peticiones se distribuyen óptimamente.
\item \textbf{Recursos del Sistema}: CPU, memoria, disco de cada nodo.
\end{itemize}

La supervisión de salud de nodos permite implementar arquitecturas auto-curativas que pueden recuperarse automáticamente de fallos.

\section{Alertas y Notificaciones}

El sistema de alertas permite configurar notificaciones para diferentes eventos:

\begin{itemize}
    \item Umbrales de uso de recursos
    \item Fallos de servicios
    \item Intentos de acceso fallidos
    \item Cambios en la configuración
\end{itemize}

\clearpage
